% ID: SEL-D0047-OA-E02
% ESTADO: cotejada. Segunda pasada visual y borrador compilado revisados.
% Fuente: Selectividad/sel_2014_matematicas_aplicadas/Matematicas_CCSS_SEPT.pdf
% SHA-256: ee5bd8e23465d55ed58339e7d9adf3d58423ac2966b53d92e39a2f9b35558ab1
% Página: 1; opción A; ejercicio original 2.

Una empresa ha realizado un estudio sobre los beneficios, en miles de euros, que ha obtenido en los últimos 10 años. La función a la que se ajustan dichos beneficios viene dada por $B(t)=2t^3-36t^2+162t-6$, con $0\leq t\leq10$.
\begin{enumerate}[label=\alph*)]
\item \textbf{(0.8 puntos)} \textquestiondown{}Qué beneficios obtuvo al inicio del periodo $(t=0)$ y al final del décimo año $(t=10)$?
\item \textbf{(1.7 puntos)} \textquestiondown{}En qué momentos se obtiene el máximo y el mínimo beneficio y cuáles fueron sus cuantías?
\end{enumerate}
